% Support + interest-to-license (combined) - Marta Parazzini (CNR-IEIIT), to Prof. Wout Joseph
% SPECTRUM 3/3 (fullest, still fully non-binding): does both in one letter - endorses the IOF
% proposal AND adds the soft "would like to use it / license it when available" hook, mirroring
% the Unisens support letter that closed with "we hope to deploy this in our projects".
% No validation, no joint paper. DRAFT to discuss together.
% \hl = "my suggestion, let's adjust", [brackets] = placeholders. Disclosure guard as before.
% Compile: pdflatex parazzini_3_support_and_license_interest.tex
\documentclass[11pt]{article}
\usepackage[a4paper,margin=2.5cm]{geometry}
\usepackage{parskip}
\usepackage{graphicx}
\usepackage{xcolor}
\definecolor{hlcol}{HTML}{B58900}
\usepackage[T1]{fontenc}
\usepackage{lmodern}
% Real yellow highlighter if soul is installed, else amber bold (wraps across lines).
\IfFileExists{soul.sty}{\usepackage{soul}\sethlcolor{yellow!60}}{}
\providecommand{\hl}[1]{\textbf{\textcolor{hlcol}{#1}}}
\pagestyle{empty}

\begin{document}

\begin{center}\small\itshape
\hl{Draft for us to go through together. Nothing here is fixed, and nothing here commits you to anything.}\\
This version does two things at once (support + a soft note of interest). We can keep both, or drop either part. Highlighted = my suggestions; [brackets] = placeholders.
\end{center}
\bigskip

% --- CNR-IEIIT letterhead placeholder ---------------------------------------
{\large\textbf{CNR-IEIIT}}\\
Istituto di Elettronica e di Ingegneria dell'Informazione e delle Telecomunicazioni\\
Consiglio Nazionale delle Ricerche, Milan, Italy\\[1.5em]
% ----------------------------------------------------------------------------

\begin{flushright}
Mr.\ Robin Wydaeghe and Prof.\ Wout Joseph\\
Ghent University -- imec\\[1em]
Milan, [date]
\end{flushright}

\bigskip
\textbf{Subject: Support and interest in the fast numerical-dosimetry tool
(\hl{AEGIS}), by Robin Wydaeghe}

\bigskip
Dear Mr.\ Wydaeghe, dear Prof.\ Joseph,

I am Director of Research at CNR-IEIIT in Milan, where my work focuses on numerical and
stochastic dosimetry of human exposure to electromagnetic fields, including \hl{5G/6G
beamforming, reconfigurable intelligent surfaces, and millimetre-wave wearable scenarios}. I
have worked together with your group and with Robin for several years, in particular within
the GOLIAT project.

Full-wave (FDTD) dosimetry is our accepted reference, but it is computationally expensive,
which in practice limits the number of scenarios, body models, and exposure conditions we can
study. The method Robin has developed is, as I understand it, very substantially faster than
FDTD while keeping comparable accuracy, and it is differentiable. \hl{A tool with these
properties would address a genuine need in our field}.

I am therefore glad to express my support for Robin Wydaeghe's IOF proposal.

\hl{And should the method become available as a software tool, I would be interested, in
principle, in using it in our dosimetry research and, if it fits our needs, in licensing it
for use in our group.} \hl{[we can soften, sharpen, or simply remove this sentence, whatever you are
comfortable with]} This note of interest is \hl{entirely non-binding} and creates no
obligation on either side.

I look forward to following how this develops.

\bigskip
With best regards,

\vspace{2.2cm}

\rule{6cm}{0.4pt}\\
Dr.\ Marta Parazzini\\
Director of Research, CNR-IEIIT\\
Consiglio Nazionale delle Ricerche, Milan, Italy\\
marta.parazzini@cnr.it

\end{document}
